\documentclass[11pt]{article}
\usepackage[margin=1in]{geometry}
\usepackage{amsmath,amssymb,amsthm}
\usepackage[T1]{fontenc}
\usepackage{lmodern}
\usepackage[colorlinks=true,linkcolor=blue,citecolor=blue,urlcolor=blue]{hyperref}
\newtheorem{theorem}{Theorem}
\newtheorem{lemma}{Lemma}
\newcommand{\BM}{\mathrm{BM}}
\newcommand{\Sset}{\mathcal S}
\newcommand{\barrier}{\mathcal B}
\title{Fourth-order diffusion approximation for three-player gambler's ruin}
\author{Kaiwen Jin\\\small\href{mailto:yui.ui.siki@gmail.com}{yui.ui.siki@gmail.com}}
\date{13 October 2026}

\begin{document}
\maketitle
\begin{abstract}
For three-player fair gambler's ruin, we prove that the probability of each
elimination order differs from its Brownian scaling limit by
$O((nN)^{-4})$ when fixed positive capitals with total $N$ are multiplied
by $n$. This is Conjecture 4.2(b) of Diaconis and Ethier. The error is antisymmetric when the two losing players are exchanged.
On the resulting half triangle, the Brownian harmonic function has a
sixth-order stencil defect with the angular factor $\sin(6\theta)$. The explicit barrier
$r^{-4}\sin(6\theta)$ and a finite discrete comparison argument give the
claimed rate.
\end{abstract}

\section{The scaling estimate}

Three players, labeled $1,2,3$, start with positive integer capitals $A,B,C$.
At each step a pair of surviving players is chosen uniformly, and a fair coin
determines a transfer of one unit between them. A player is removed on reaching
zero. Write $P_{A,B,C}(\sigma)$ for the probability of elimination order
$\sigma\in S_3$, with the winner listed last, and set $N=A+B+C$.
Diaconis and Ethier \cite[Theorem 4.1]{DE} identify the limits
\[
 P^{\BM}_{A,B,C}(\sigma)=\lim_{m\to\infty}P_{mA,mB,mC}(\sigma)
\]
by Brownian exit distributions in an equilateral triangle. Their
Conjecture 4.2(b) asks for a fourth-order error bound.
Denisov and Wachtel \cite[Proposition 5]{DW} prove a cubic bound and suggest
that symmetry may yield the faster rate. Their asymptotics with two fixed
capitals and a growing third capital concern a different regime.

\begin{theorem}\label{thm:main}
For every fixed positive integer triple $A,B,C$ and every $\sigma\in S_3$,
\[
 \left|P_{nA,nB,nC}(\sigma)-P^{\BM}_{A,B,C}(\sigma)\right|
 =O\bigl((nN)^{-4}\bigr),\qquad n\to\infty.
\]
The implied constant may depend on the capital proportions and on $\sigma$.
\end{theorem}

The sixth-order local defect alone does not explain this rate: the boundary
payoff jumps at one corner. Exchanging the two losing players makes the error
odd across the bisector of that corner. This reduces the comparison to a
$30$-degree wedge, where a positive barrier of degree $-4$ is available.
The lattice has one jump across the bisector; retaining its exact contribution
is necessary for the comparison.

\section{Harmonic functions and reflection}

Put $\omega=e^{i\pi/3}$ and use the lattice and triangle
\[
 \Lambda=\{b+c\omega:b,c\in\mathbb Z\},\qquad
 T_M=\{b+c\omega:b,c>0,\ b+c<M\}.
\]
Here $M$ is an integer total capital, $a=M-b-c$, and $z=b+c\omega$.
The transition before the first elimination has the six equally likely steps
\[
 \Sset=\{\pm1,\ \pm\omega,\ \pm(\omega-1)\}.
\]
For functions defined at the necessary neighbors, let
\[
 Ph(z)=\frac16\sum_{s\in\Sset}h(z+s),\qquad L=I-P.
\]
Define $v_M(z)=P_{a,b,c}(321)$. The remaining two-player game has winning
probability equal to the player's fraction of the total capital: that affine
function satisfies the averaging equation and the two absorbing values.
For the two-player chain, absorption is almost sure: a prescribed sequence
of at most $M$ transfers in one direction has probability at least $2^{-M}$
and reaches an absorbing endpoint. In the three-player chain, a prescribed
sequence of at most $M$ transfers for a fixed pair and direction has probability
at least $6^{-M}$ and forces an elimination. Repeating finite blocks proves
almost-sure elimination, followed by the two-player absorption just established. Consequently the boundary data for
$v_M$ are
\begin{equation}\label{eq:boundary}
 g_M(z)=1-b/M\quad(c=0),\qquad
 g_M(z)=0\quad(b=0\text{ or }b+c=M).
\end{equation}
These prescriptions refer to open sides. An interior lattice point cannot
jump directly to a vertex. The finite averaging problem is unique:
a positive maximum of a zero-boundary difference propagates along the steps
$(b,c)\mapsto(b,c-1)$ until it reaches the boundary. Thus $Lv_M=0$ and
$0\le v_M\le1$.

Let $u_M$ be the bounded harmonic function with the same open-side data,
given by the Brownian exit expectation. It satisfies $0\le u_M\le1$.
The map $q=2z/M-1$ sends our triangle to the one used in
\cite[Theorem 4.1]{DE}. On its bottom side the weight $|q-1|/2$ is $1-b/M$.
Similarities preserve isotropic Brownian exit distributions, so
\begin{equation}\label{eq:bm}
 u_M(z)=u_1(z/M),\qquad
 u_{nN}\bigl(n(B+C\omega)\bigr)=P^{\BM}_{A,B,C}(321).
\end{equation}

We shall use odd reflection for zero-trace harmonic functions, removal of
bounded isolated harmonic singularities, and harmonic derivative estimates.
These are, respectively, Theorems 4.12, 2.3 and 2.4 of
Axler, Bourdon and Ramey \cite{HFT}.
At a corner of angle $\pi/3$, six successive odd reflections give a bounded
harmonic function on a punctured disk whenever both adjacent traces are zero.
The sectors agree on their zero-trace rays and are consistent after a full
turn. The singularity is removable, and oddness forces the value at the
corner to be zero. It follows by the maximum principle
\cite[Corollary 1.10]{HFT} that bounded harmonic functions with the same
open-side data on $T_M$ are unique.

\begin{lemma}\label{lem:symmetry}
Let $Rz=\omega\overline z$, which interchanges $b$ and $c$, and put
$\ell_M(z)=1-(b+c)/M$. Then
\begin{equation}\label{eq:symmetry}
 v_M(z)+v_M(Rz)=\ell_M(z),\qquad
 u_M(z)+u_M(Rz)=\ell_M(z).
\end{equation}
In particular $e_M=v_M-u_M$ is odd under $R$ and is zero on $b=c$.
\end{lemma}
\begin{proof}
The reflection preserves $\Sset$ and $T_M$. On each open side the sum of the
two boundary values in \eqref{eq:boundary} is $\ell_M$.
This affine function is harmonic and satisfies $P\ell_M=\ell_M$.
Discrete and continuous uniqueness prove \eqref{eq:symmetry}.
\end{proof}

Only the corner $0$ has incompatible boundary traces. At the other two
corners we can reflect after subtracting an affine function matching both
adjacent sides. This gives derivative bounds which do not deteriorate near
those corners.

\begin{lemma}\label{lem:derivatives}
There are constants $\kappa,C_0>0$, independent of $M$, such that for every
$z\in\overline T_M\setminus\{0\}$, $u_M$ has a harmonic continuation to
$B(z,\kappa|z|)$ bounded in absolute value by $C_0$.
For every nonnegative integer $k$ this implies
\begin{equation}\label{eq:derivatives}
 |D^\alpha u_M(z)|\le C_k|z|^{-k},\qquad |\alpha|=k,
\end{equation}
with $C_k$ independent of $M$.
\end{lemma}
\begin{proof}
It suffices to work with $M=1$. At an open side, subtract an affine harmonic
function with the prescribed trace, reflect the difference oddly, and add the
affine function back. At vertex $1$, $\ell_1$ matches both adjacent traces.
At vertex $\omega$, both traces are zero. The six-reflection argument above
therefore gives harmonic continuations through these vertices.

Near $0$, a ball of sufficiently small fixed relative radius $\kappa|z|$
misses $0$ and the opposite side, and meets at most one of the two adjacent
side lines. A single affine reflection then gives a uniformly bounded
continuation. Away from $0$, the remaining closed part of the normalized
triangle is compact. A finite cover by the side and vertex continuations,
followed by a reduction of $\kappa$, gives the asserted balls and a uniform
bound. The continuations agree on overlaps by real analyticity
\cite[Theorems 1.27--1.28]{HFT}; overlapping local disks can be chosen to meet
the original triangle. The harmonic derivative estimate on these balls gives
\eqref{eq:derivatives}. Scaling by $M$ proves both conclusions.
\end{proof}

\section{The defect and a barrier}

For a uniformly chosen step $s\in\Sset$, the odd moments vanish and
\[
 \mathbb E(s_x^2)=\mathbb E(s_y^2)=\tfrac12,\quad
 \mathbb E(s_xs_y)=0,
 \quad \mathbb E(s_x^4)=\mathbb E(s_y^4)=\tfrac38,
 \quad \mathbb E(s_x^2s_y^2)=\tfrac18,
\]
while $\mathbb E(s_x^3s_y)=\mathbb E(s_xs_y^3)=0$.
Taylor expansion through degree five therefore yields
\begin{equation}\label{eq:taylor}
 Ph-h=\tfrac14\Delta h+\tfrac1{64}\Delta^2h+\mathcal R_h,
 \qquad |\mathcal R_h(z)|\le C\sup_{|w-z|\le1}|D^6h(w)|.
\end{equation}
Here $|D^k h|$ denotes the maximum absolute value of its partial derivatives
of total order $k$. One may obtain the remainder directly by applying
\[
 h(z+s)=\sum_{j=0}^5\frac{(s\cdot\nabla)^j h(z)}{j!}
 +\frac1{5!}\int_0^1(1-t)^5(s\cdot\nabla)^6h(z+ts)\,dt
\]
to each step. Applying the same formula to the first derivatives also gives
\[
 |\nabla\mathcal R_h(z)|\le C\sup_{|w-z|\le1}|D^7h(w)|.
\]

Write $z=re^{i\theta}$. The half triangle $b>c>0$ lies in
$0<\theta<\pi/6$.

\begin{lemma}\label{lem:defect}
There are fixed constants $r_0,C_f>0$, independent of $M$, such that
$f_M=Pu_M-u_M$ obeys
\begin{equation}\label{eq:defect}
 |f_M(z)|\le C_f r^{-6}\sin(6\theta)
 \quad(z\in T_M,\ b>c,\ r\ge r_0).
\end{equation}
The stencil here uses the harmonic continuations when $z$ is not a lattice
point and a shifted point lies outside $T_M$.
\end{lemma}
\begin{proof}
Choose $r_0$ large enough that the relative balls of
Lemma~\ref{lem:derivatives} contain all the unit neighborhoods under
consideration. Harmonicity and \eqref{eq:taylor} give
\begin{equation}\label{eq:rawdefect}
 |f_M(z)|\le Cr^{-6},\qquad |\nabla f_M(z)|\le Cr^{-7}.
\end{equation}
The same bounds hold along segments whose radii are a fixed fraction of $r$.

On the lower side $c=0$, subtract $1-\operatorname{Re}z/M$ before odd
reflection. Since $P$ is exact on affine functions and commutes with complex
conjugation, $f_M$ is odd about that side and is zero there.
Likewise \eqref{eq:symmetry} and invariance under $R$ make $f_M$ zero on
the bisector. These identities hold for the local continuations too, by
analyticity, including when a stencil near the bisector meets the opposite
side of the triangle.

Project $z$ perpendicularly onto the nearer ray $\theta=0$ or
$\theta=\pi/6$. The lower-side foot has capital coordinates $(b+c/2,0)$;
the bisector foot has coordinates $((b+c)/2,(b+c)/2)$.
Both feet and the connecting segments remain in the closed triangle.
Every radius on the chosen segment is at least $r\cos(\pi/12)$.
Set $t=\min(\theta,\pi/6-\theta)$. The segment length is $r\sin t$.
The zero at the foot and the gradient estimate in \eqref{eq:rawdefect} give
$|f_M(z)|\le Cr^{-6}\sin t$. On $0\le t\le\pi/12$,
$\sin t\le C\sin(6t)$, and $\sin(6t)=\sin(6\theta)$.
This proves \eqref{eq:defect}.
\end{proof}

\begin{lemma}\label{lem:barrier}
The function
\[
 \barrier(z)=\frac{\operatorname{Im}(z^6)}{|z|^{10}}
           =r^{-4}\sin(6\theta)
\]
is positive for $0<\theta<\pi/6$. After increasing a fixed $r_0$ if necessary,
\begin{equation}\label{eq:barrier}
 L\barrier(z)\ge\tfrac52 r^{-6}\sin(6\theta),\qquad r\ge r_0,
 \quad 0<\theta<\pi/6.
\end{equation}
\end{lemma}
\begin{proof}
The polar-coordinate identity
$\Delta(r^q\sin(k\theta))=(q^2-k^2)r^{q-2}\sin(k\theta)$ gives
\[
 \Delta\barrier=-20r^{-6}\sin(6\theta),\qquad
 \Delta^2\barrier=0.
\]
By smoothness on the unit circle and homogeneity,
$|D^j\barrier(z)|\le C_jr^{-4-j}$. Thus \eqref{eq:taylor} gives
\[
 L\barrier=5r^{-6}\sin(6\theta)+\rho,
 \qquad |\rho|\le Cr^{-10},\quad |\nabla\rho|\le Cr^{-11}.
\]
Both $\barrier$ and the leading term are odd about the two bounding rays,
and $L$ commutes with their reflections. Hence $\rho$ is zero on both rays.
The projection argument in Lemma~\ref{lem:defect} now gives
$|\rho|\le Cr^{-10}\sin(6\theta)$. Choose the fixed $r_0$ so that
$Cr_0^{-4}\le5/2$ to obtain \eqref{eq:barrier}.
\end{proof}

\section{Discrete comparison and the full estimate}

A step from $b>c$ can cross to $b'<c'$ only if $b-c=1$.
There is exactly one such step, namely
$(b,c)\mapsto(b-1,c+1)=(c,b)$. For an odd function its value at that
neighbor is minus its value at the original point. Thus, when restricted to
the half triangle, $L$ has diagonal
$1+\tfrac16\mathbf1_{\{b-c=1\}}$ and coefficient $-1/6$ at every remaining
neighbor; fixed-line values are zero.

This operator has the finite comparison property. Indeed, suppose $Lw\le0$
on a finite subset of the half triangle and $w\le0$ at its outside
neighbors, with the crossing value interpreted by odd reflection.
If a positive maximum $m$ occurs where $b-c=1$, then
$Lw\ge(7/6)m-(5/6)m=m/3>0$. At any other point, $Lw\ge0$,
and equality forces all six neighbor values to equal $m$.
Following $(b,c)\mapsto(b,c-1)$ then reaches an outside neighbor or $c=0$,
contradicting its nonpositive value. Hence $w\le0$ on the subset.

\begin{proof}[Proof of Theorem~\ref{thm:main}]
Since $Lv_M=0$, the error satisfies $Le_M=f_M$.
Use the same fixed $r_0$ in Lemmas~\ref{lem:defect} and \ref{lem:barrier}, and set
\[
 E_M=\{z=b+c\omega\in T_M\cap\Lambda:b>c,\ |z|\ge r_0\}.
\]
Choose $K\ge2C_f/5$. Increase it, independently of $M$, so that
$K\barrier(z)\ge1$ at all half-wedge lattice points with $0<|z|<r_0$.
There are only finitely many such points, and $\barrier$ is positive at each.
For either sign,
\[
 L(\pm e_M-K\barrier)\le0\quad\hbox{on }E_M.
\]
At inner neighbors outside $E_M$, $|e_M|\le1\le K\barrier$.
At physical sides $e_M=0$ and $\barrier\ge0$; on the bisector both are zero.
There is no jump over a physical side or directly to a vertex, and the sole
crossing jump is already included in the half-domain diagonal.
Finite comparison therefore yields
\begin{equation}\label{eq:halferror}
 |e_M(z)|\le K|z|^{-4}\sin(6\theta)\quad(z\in E_M).
\end{equation}
Reflection by $R$, and the exact zero on its fixed line, give
\begin{equation}\label{eq:radial}
 |v_M(z)-u_M(z)|\le K|z|^{-4}
 \quad(z\in T_M\cap\Lambda,\ |z|\ge r_0),
\end{equation}
with the same constants for every $M$.

For fixed $A,B,C$, take $M=nN$ and $z=n(B+C\omega)$ in
\eqref{eq:radial}. Since $|z|^2=n^2(B^2+BC+C^2)$, \eqref{eq:bm} gives,
for all sufficiently large integers $n$,
\[
 \left|P_{nA,nB,nC}(321)-P^{\BM}_{A,B,C}(321)\right|
 \le\frac{K}{n^4(B^2+BC+C^2)^2}
 =\frac{KN^4}{(B^2+BC+C^2)^2}(nN)^{-4}.
\]
The last coefficient depends only on the fixed proportions. Relabeling the
players preserves the transition law and proves the estimate for every
$\sigma\in S_3$. When the two losing players start with equal capitals,
\eqref{eq:symmetry} gives exact agreement with the Brownian probability.
\end{proof}

\paragraph{AI assistance.}
OpenAI Codex assisted with the argument, proof drafting and algebraic checks.
This is a natural-language proof submission. No Lean verification or human
referee certification is asserted. The submission records the scope of the
AI review performed.

\begin{thebibliography}{9}
\bibitem{HFT}
S. Axler, P. Bourdon and W. Ramey,
\emph{Harmonic Function Theory}, 2nd ed., Graduate Texts in Mathematics 137,
Springer, 2001.
\href{https://www.axler.net/HFT.pdf}{Author-hosted book}.
\bibitem{DE}
P. Diaconis and S. N. Ethier, Gambler's ruin and the ICM,
\emph{Statistical Science} \textbf{37}(3) (2022), 289--305.
\href{https://doi.org/10.1214/21-STS826}{doi:10.1214/21-STS826}.
Author manuscript: \href{https://arxiv.org/abs/2011.07610v3}{arXiv:2011.07610v3}.
\bibitem{DW}
D. Denisov and V. Wachtel, Harmonic measure in a multidimensional gambler's problem,
\emph{Annals of Applied Probability} \textbf{34}(5) (2024), 4387--4407.
\href{https://doi.org/10.1214/24-AAP2069}{doi:10.1214/24-AAP2069}.
Author manuscript: \href{https://arxiv.org/abs/2212.11526v1}{arXiv:2212.11526v1}.
\end{thebibliography}
\end{document}
